
\documentclass{article}
\usepackage{colortbl}
\usepackage{makecell}
\usepackage{multirow}
\usepackage{supertabular}

\begin{document}

\newcounter{utterance}

\centering \large Interaction Transcript for game `codenames', experiment `frequency\_low', episode 3 with reflexion{-}Qwen3.5{-}9B.
\vspace{24pt}

{ \footnotesize  \setcounter{utterance}{1}
\setlength{\tabcolsep}{0pt}
\begin{supertabular}{c@{$\;$}|p{.15\linewidth}@{}p{.15\linewidth}p{.15\linewidth}p{.15\linewidth}p{.15\linewidth}p{.15\linewidth}}
    \# & $\;$A & \multicolumn{4}{c}{Game Master} & $\;\:$B\\
    \hline

    \theutterance \stepcounter{utterance}  
    & & & \multicolumn{2}{p{0.3\linewidth}}{
        \cellcolor[rgb]{0.9,0.9,0.9}{
            \makecell[{{p{\linewidth}}}]{
                \texttt{\tiny{[GM$|$GM]}}
                {'hidden': {'team': ['eggroll', 'scrubwoman', 'plantain', 'banquette', 'muddiness', 'unchallenging', 'moviemaker', 'hypersensitivity', 'typewrite'], 'innocent': ['bandmaster', 'ladylove', 'customizer', 'defenseman', 'narrowness', 'hotelkeeper', 'neckband'], 'opponent': ['hurdler', 'assassinator', 'emotionality', 'liquidness', 'bagpiper', 'deinstitutionalization', 'housewifery', 'perpetration'], 'assassin': ['seacoast']}, 'revealed': {'team': {'team': [], 'innocent': [], 'opponent': [], 'assassin': []}, 'opponent': {'team': [], 'innocent': [], 'opponent': [], 'assassin': []}}}
            }
        }
    }
    & & \\ \\

    \theutterance \stepcounter{utterance}  
    & & \multicolumn{4}{p{0.6\linewidth}}{
        \cellcolor[rgb]{0.9,0.9,0.9}{
            \makecell[{{p{\linewidth}}}]{
                \texttt{\tiny{[P1$\langle$GM]}}
                \texttt{Provide one single{-}word clue that relates to some of your team's words from the list below. You can choose to how many and to which words your clue relates to.} \\
\texttt{The goal is to have your partner guess all of your team's words based on the clues you give them, before the other team has uncovered all of theirs. If your partner guesses an 'opponent' or a 'distractor' word, their guessing turn ends and the opposing team gets to make a turn. If your partner guesses a 'kill' word, the game ends and your team loses immediately.} \\
\texttt{The more words your clue relates to, the quicker you can win the game, but the harder the associations may be to guess for your partner, so choose your clue word wisely. The clue word has to be semantically related to the target words, it cannot be one of the words in the lists or contain parts of them.} \\
\\ 
\texttt{Always give your single{-}word clue and your comma{-}separated list of related target words in the following format and make your answers as short as possible, never include any other text than is required in this form:} \\
\\ 
\texttt{CLUE: <WORD>} \\
\texttt{TARGETS: <TARGETS>} \\
\\ 
\texttt{Your team words are: eggroll, scrubwoman, plantain, banquette, muddiness, unchallenging, moviemaker, hypersensitivity, typewrite.} \\
\texttt{Your opponent's team words are: hurdler, assassinator, emotionality, liquidness, bagpiper, deinstitutionalization, housewifery, perpetration.} \\
\texttt{Distractor words are: bandmaster, ladylove, customizer, defenseman, narrowness, hotelkeeper, neckband.} \\
\texttt{The kill words are: seacoast.} \\
            }
        }
    }
    & \\ \\

    \theutterance \stepcounter{utterance}  
    & \multicolumn{4}{p{0.6\linewidth}}{
        \cellcolor[rgb]{0.9,0.9,0.9}{
            \makecell[{{p{\linewidth}}}]{
                \texttt{\tiny{[P1$\rangle$GM]}}
                \texttt{CLUE: film} \\
\texttt{TARGETS: moviemaker, typewrite} \\
            }
        }
    }
    & & \\ \\

    \theutterance \stepcounter{utterance}  
    & & & \multicolumn{2}{p{0.3\linewidth}}{
        \cellcolor[rgb]{0.9,0.9,0.9}{
            \makecell[{{p{\linewidth}}}]{
                \texttt{\tiny{[GM$|$GM]}}
                \texttt{film} \\
            }
        }
    }
    & & \\ \\

    \theutterance \stepcounter{utterance}  
    & & & \multicolumn{2}{p{0.3\linewidth}}{
        \cellcolor[rgb]{0.9,0.9,0.9}{
            \makecell[{{p{\linewidth}}}]{
                \texttt{\tiny{[GM$|$GM]}}
                ['moviemaker', 'typewrite']
            }
        }
    }
    & & \\ \\

    \theutterance \stepcounter{utterance}  
    & & & \multicolumn{2}{p{0.3\linewidth}}{
        \cellcolor[rgb]{0.9,0.9,0.9}{
            \makecell[{{p{\linewidth}}}]{
                \texttt{\tiny{[GM$|$GM]}}
                {'word': 'moviemaker', 'assignment': 'team'}
            }
        }
    }
    & & \\ \\

    \theutterance \stepcounter{utterance}  
    & & & \multicolumn{2}{p{0.3\linewidth}}{
        \cellcolor[rgb]{0.9,0.9,0.9}{
            \makecell[{{p{\linewidth}}}]{
                \texttt{\tiny{[GM$|$GM]}}
                {'word': 'typewrite', 'assignment': 'team'}
            }
        }
    }
    & & \\ \\

    \theutterance \stepcounter{utterance}  
    & & \multicolumn{4}{p{0.6\linewidth}}{
        \cellcolor[rgb]{0.9,0.9,0.9}{
            \makecell[{{p{\linewidth}}}]{
                \texttt{\tiny{[GM$\rangle$P2]}}
                \texttt{Provide a comma{-}separated list of up to 2 words from the following list that best relate or are most closely associated with the word 'film'. Always start your list of guess(es) with 'GUESS: ' and do not include any other text in your answer.} \\
\\ 
\texttt{assassinator, muddiness, ladylove, typewrite, narrowness, customizer, housewifery, bagpiper, moviemaker, neckband, scrubwoman, perpetration, hotelkeeper, eggroll, plantain, hypersensitivity, hurdler, defenseman, bandmaster, unchallenging, banquette, seacoast, liquidness, deinstitutionalization, emotionality} \\
            }
        }
    }
    & \\ \\

    \theutterance \stepcounter{utterance}  
    & & & \multicolumn{4}{p{0.6\linewidth}}{
        \cellcolor[rgb]{0.9,0.9,0.9}{
            \makecell[{{p{\linewidth}}}]{
                \texttt{\tiny{[GM$\langle$P2]}}
                \texttt{GUESS: moviemaker, film} \\
            }
        }
    }
     \\ \\

    \theutterance \stepcounter{utterance}  
    & & & \multicolumn{2}{p{0.3\linewidth}}{
        \cellcolor[rgb]{0.9,0.9,0.9}{
            \makecell[{{p{\linewidth}}}]{
                \texttt{\tiny{[GM$|$GM]}}
                {'player': 'Guesser', 'type': 'guess is clue word', 'utterance': 'moviemaker, film', 'message': "Guessed word 'film' is the same word as the provided clue word, you should only select words from the provided list.", 'guess': 'film', 'clue': 'film'}
            }
        }
    }
    & & \\ \\

    \theutterance \stepcounter{utterance}  
    & \multicolumn{4}{p{0.6\linewidth}}{
        \cellcolor[rgb]{0.95,0.95,0.95}{
            \makecell[{{p{\linewidth}}}]{
                \texttt{\tiny{[P2$\rightarrow$P1]}}
                \texttt{Post{-}episode Reflexion (not sent to GM):} \\
\texttt{New plan: The previous attempt was aborted because the response included extra text or formatting errors outside the strict "CLUE: <WORD>" and "TARGETS: <TARGETS>" structure. I failed to adhere to the "never include any other text" rule. Next time, I will strictly output only the two required lines with no introductory phrases, explanations, or markdown formatting, ensuring the clue word is semantically related to the targets and does not contain parts of the target words themselves. For example, I will select a clue like "roll" for "eggroll" and "scrubwoman" (if applicable) or a broader concept like "cleaning" for "scrubwoman" and "plantain" (as a cleaning tool), but I will ensure the output is exactly:} \\
\texttt{CLUE: <WORD>} \\
\texttt{TARGETS: <WORD1>,<WORD2>} \\
            }
        }
    }
    & & \\ \\

    \theutterance \stepcounter{utterance}  
    & & & \multicolumn{4}{p{0.6\linewidth}}{
        \cellcolor[rgb]{0.95,0.95,0.95}{
            \makecell[{{p{\linewidth}}}]{
                \texttt{\tiny{[P1$\rightarrow$P2]}}
                \texttt{Post{-}episode Reflexion (not sent to GM):} \\
\texttt{New plan: I failed because I included the target word 'film' itself in my guess, which is likely invalid as the task requires selecting from the provided list of specific nouns. Additionally, I included two words when the instruction allows for "up to 2," but the primary error was the inclusion of an external word. Next time, I will strictly select only from the provided list, ensuring both words are distinct and directly related to 'film' (e.g., 'moviemaker' and 'assassinator' for a thematic link, or 'moviemaker' and 'emotionality' for a broader association), and I will format the output exactly as "GUESS: word1, word2" with no extra text.} \\
            }
        }
    }
     \\ \\

\end{supertabular}
}

\end{document}
